\documentclass[conference]{IEEEtran}

\usepackage{booktabs}
\usepackage{graphicx}
\usepackage{xcolor}
\usepackage{hyperref}
\usepackage{amsmath}
% algorithm packages removed - using manual pseudocode instead
\usepackage{multirow}
% enumitem not available in TinyTeX

\hypersetup{colorlinks=true, linkcolor=blue, citecolor=blue, urlcolor=blue}

\begin{document}

\title{SADE: Symptom-Aware Diagnostic Escalation\\for LLM-Based Network Fault Diagnosis}

\author{
\IEEEauthorblockN{[Author Name]}
\IEEEauthorblockA{[Affiliation]\\
[Email]}
}

\maketitle

\begin{abstract}
Automated network fault diagnosis using Large Language Models (LLMs) has emerged as a promising approach to reduce manual troubleshooting effort. However, existing LLM-based agents employ flat, unstructured exploration strategies that scale poorly with network complexity. We present SADE (Symptom-Aware Diagnostic Escalation), a structured diagnostic agent architecture that combines layered symptom-driven investigation (L1: configuration, L2: infrastructure, L3: behavioral) with modular, on-demand diagnostic skills. We evaluate SADE on the NIKA network troubleshooting benchmark across 36 fault types and three topology sizes. SADE achieves an RCA accuracy of \textbf{XX\%}, compared to 44.8\% for the baseline ReAct/GPT-5 agent and 67.6\% for an unstructured Claude agent, demonstrating that domain-structured prompting significantly outperforms generic exploration for network diagnosis. We additionally identify and characterize a class of \emph{self-healing faults} in container-based network emulation environments where fault evidence disappears before any agent can observe it, contributing a benchmark design finding for future evaluation work.
\end{abstract}

\begin{IEEEkeywords}
network fault diagnosis, large language models, automated troubleshooting, benchmark evaluation, prompt engineering
\end{IEEEkeywords}

% ============================================================
\section{Introduction}
% ============================================================

Network fault diagnosis remains one of the most time-consuming tasks in network operations. A single fault---such as a misconfigured firewall rule, a poisoned ARP cache, or a resource-exhausted server---can manifest as degraded service across dozens of hosts, requiring operators to systematically narrow down the root cause from network-wide symptoms.

Recent advances in Large Language Models (LLMs) have enabled autonomous agents that can interact with network devices through tool calls, interpret diagnostic outputs, and reason about fault causality~\cite{nika}. The NIKA benchmark~\cite{nika} provides a standardized evaluation framework for such agents, using Kathara-based container emulation to inject 36 distinct fault types across varying topology sizes.

However, existing LLM-based diagnostic agents typically employ a flat ReAct~\cite{react} strategy: observe symptoms, think about what to check next, act, repeat. This approach has two critical limitations. First, the agent lacks a systematic prioritization of diagnostic layers---it may exhaust tool-call budgets investigating application-layer symptoms when the root cause is a Layer~2 MAC conflict. Second, the agent carries all diagnostic knowledge in a single monolithic prompt, leading to token inefficiency and difficulty maintaining fault-specific expertise.

We propose \textbf{SADE} (Symptom-Aware Diagnostic Escalation), an agent architecture that addresses both limitations:

\begin{itemize}    \item \textbf{Layered diagnostic workflow.} SADE follows a deterministic escalation path---configuration sweep (L1), infrastructure and ACL sweep (L2), behavioral analysis (L3)---ensuring systematic coverage before concluding.
    \item \textbf{Modular diagnostic skills.} Instead of embedding all fault fingerprints in the system prompt, SADE loads fault-specific diagnostic skills on-demand based on observed symptoms, reducing prompt size while maintaining deep expertise.
    \item \textbf{Self-healing fault characterization.} Through systematic failure analysis, we identify fault types in container-based emulation where evidence disappears before observation, contributing a benchmark design finding.
\end{itemize}

We evaluate SADE on the NIKA benchmark and compare against two baselines: the original ReAct/GPT-5 agent and an unstructured Claude-based agent. Our results show that SADE substantially improves root cause analysis accuracy, particularly for fault types requiring cross-layer correlation.

% ============================================================
\section{Related Work}
% ============================================================

\textbf{LLM agents for network operations.}
The application of LLMs to network management has gained significant attention. Prior work has explored using LLMs for configuration generation, anomaly detection, and intent-based networking. However, autonomous diagnostic agents that interact directly with network devices remain underexplored.

\textbf{Tool-augmented reasoning.}
ReAct~\cite{react} introduced interleaved reasoning and acting for LLM agents. Subsequent work has extended this with structured planning, skill libraries, and hierarchical task decomposition. SADE builds on these ideas by introducing domain-specific diagnostic escalation.

\textbf{Network fault diagnosis benchmarks.}
NIKA~\cite{nika} provides a comprehensive benchmark for evaluating LLM-based network troubleshooting agents, with 36 fault types across container-emulated topologies. Our work evaluates on this benchmark and additionally characterizes limitations in its fault injection mechanism.

% ============================================================
\section{Background: NIKA Benchmark}
% ============================================================

NIKA~\cite{nika} is a network troubleshooting benchmark built on Kathara, a container-based network emulation platform. Each evaluation scenario consists of:

\begin{itemize}    \item A network topology (\texttt{ospf\_enterprise\_dhcp}) with hosts, routers, switches, web servers, a load balancer, and DNS/DHCP servers.
    \item A single injected fault from one of 36 types spanning six categories: end-host failures, ACL blocks, link faults, traffic control issues, resource contention, and routing problems.
    \item Three topology sizes: small (\texttt{s}, $\sim$10 hosts), medium (\texttt{m}, $\sim$30 hosts), and large (\texttt{l}, $\sim$130 hosts).
\end{itemize}

The agent interacts with the network through MCP (Model Context Protocol) tools: \texttt{exec\_shell} for running commands on devices, \texttt{get\_reachability} for ping-based connectivity tests, \texttt{curl\_web\_test} for HTTP response timing, and \texttt{get\_tc\_statistics} for traffic control state. The agent must submit a diagnosis containing: (1) whether an anomaly exists, (2) which devices are faulty, and (3) the root cause name.

Evaluation uses three metrics: \textbf{detection accuracy} (is the network anomalous?), \textbf{localization accuracy} (which devices are faulty?), and \textbf{RCA accuracy} (what is the root cause?).

% ============================================================
\section{SADE Agent Architecture}
\label{sec:sade}
% ============================================================

\subsection{Design Principles}

SADE is motivated by two observations from analyzing baseline agent failures:

\begin{enumerate}    \item \textbf{Layer confusion.} Flat ReAct agents frequently misattribute symptoms across network layers. For example, an agent observing HTTP timeouts may conclude ``server overload'' when the actual root cause is an ARP cache poisoning at Layer~2 that prevents outbound traffic.
    \item \textbf{Prompt saturation.} Embedding diagnostic fingerprints for all 36 fault types in a single prompt leads to token waste and diluted attention over fault-specific details.
\end{enumerate}

SADE addresses these through a layered workflow and modular skill system.

\subsection{Layered Diagnostic Workflow}

The following box shows the SADE diagnostic workflow. The key insight is that each layer acts as a \emph{filter}: if a fault is detected at a lower layer, the agent reads the corresponding skill and submits immediately. Only when a layer is clean does the agent escalate to the next.

\vspace{4pt}
\noindent\fbox{\parbox{0.96\columnwidth}{
\small
\textbf{SADE Diagnostic Workflow} \\[2pt]
\textbf{Step 0a:} Load topology info + baseline behavior \\
\textbf{Step 0b:} Host config sweep + ARP check \\
\hspace*{1em} \textit{if} host config differs $\rightarrow$ load \texttt{host-ip-conflict-skill}, submit \\
\hspace*{1em} \textit{if} ARP MAC mismatch $\rightarrow$ load \texttt{host-ip-conflict-skill}, submit \\
\textbf{Step 0c:} Infrastructure sweep (ACL, link state, MAC) \\
\hspace*{1em} \textit{if} nft DROP rule $\rightarrow$ load \texttt{iptables-acl-skill}, submit \\
\hspace*{1em} \textit{if} interface DOWN/ABSENT $\rightarrow$ load \texttt{link-fault-skill}, submit \\
\hspace*{1em} \textit{if} duplicate MACs $\rightarrow$ load \texttt{host-ip-conflict-skill}, submit \\
\textbf{Step 1:} Reachability test (L3 connectivity) \\
\textbf{Step 2:} HTTP response timing + TC statistics \\
\hspace*{1em} \textit{if} TC qdisc anomaly $\rightarrow$ load \texttt{tc-fault-skill}, submit \\
\textbf{Step 3:} Resource + application code check \\
\hspace*{1em} Check \texttt{ps aux}, web server source $\rightarrow$ load \texttt{resource-contention-skill} \\
\textbf{Stall Recovery:} exhaustive re-check if all layers clean
}}\vspace{4pt}

The workflow is encoded in the agent's system prompt as a concise skeleton ($\sim$60 lines), with each step containing trigger conditions that point to the relevant skill file. This contrasts with the baseline approach of embedding all diagnostic logic inline.

\subsection{Modular Diagnostic Skills}

SADE uses 11 diagnostic skill files, each containing fault-specific fingerprints, disambiguation rules, and submission guidelines. Skills are stored as Markdown files in the agent's working directory and loaded via file reads when triggered by workflow conditions.

Table~\ref{tab:skills} lists the skill modules. Each skill is self-contained: it describes how to detect faults in its category, how to distinguish between similar fault types, and what to submit. This modular design has three benefits:

\begin{table}[t]
\centering
\caption{SADE Diagnostic Skill Modules}
\label{tab:skills}
\small
\begin{tabular}{@{}ll@{}}
\toprule
\textbf{Skill Module} & \textbf{Fault Types Covered} \\
\midrule
\texttt{iptables-acl-skill} & ARP/ICMP/HTTP/OSPF/BGP ACL blocks \\
\texttt{link-fault-skill} & link\_flap, link\_down, link\_detach \\
\texttt{tc-fault-skill} & bandwidth throttling, incast, fragmentation \\
\texttt{host-ip-conflict-skill} & IP/MAC conflict, ARP poisoning, misconfig \\
\texttt{dns-fault-skill} & incorrect DNS, record error \\
\texttt{bgp-fault-skill} & BGP hijack, blackhole route leak \\
\texttt{load-balancer-skill} & load balancer overload \\
\texttt{resource-contention-skill} & sender/receiver contention, app delay, DoS \\
\texttt{baseline-behavior-skill} & normal behavior reference (non-faults) \\
\texttt{diagnosis-methodology-skill} & general diagnostic strategies \\
\texttt{big-return-skill} & large output parsing utilities \\
\bottomrule
\end{tabular}
\end{table}

\begin{enumerate}    \item \textbf{Token efficiency.} The system prompt remains compact. Skills are loaded only when relevant, avoiding wasted context on unrelated fault types.
    \item \textbf{Maintainability.} Each skill can be independently updated without affecting the core workflow or other skills.
    \item \textbf{Targeted attention.} When a skill is loaded, the agent's full attention is focused on the relevant fault category, reducing confusion between similar fault types.
\end{enumerate}

\subsection{Symptom-to-Skill Routing}

The mapping from observed symptoms to skills is embedded in the workflow steps as brief trigger conditions. For example:

\begin{itemize}    \item \texttt{nft list ruleset} shows DROP $\rightarrow$ load \texttt{iptables-acl-skill}
    \item \texttt{ip link show} shows interface ABSENT $\rightarrow$ load \texttt{link-fault-skill}
    \item One web server's source code differs $\rightarrow$ load \texttt{resource-contention-skill}
\end{itemize}

This routing is deterministic and exhaustive: the workflow guarantees that all fault categories are checked, and the stall recovery phase provides a final sweep for edge cases.

% ============================================================
\section{Experimental Setup}
% ============================================================

\subsection{Agents Compared}

We evaluate three agent configurations:

\begin{enumerate}    \item \textbf{ReAct/GPT-5} (Baseline): The NIKA default agent using a ReAct prompting strategy with GPT-5 as the backend LLM.
    \item \textbf{Claude-base}: Claude Sonnet~4 with a basic diagnostic prompt (no layered workflow, no skills).
    \item \textbf{SADE}: Claude Sonnet~4 with the SADE layered workflow and 11 diagnostic skills.
\end{enumerate}

All agents use the same MCP tool interface and are subject to the same turn budget (max 20 turns). The only differences are the system prompt, available skills, and backend LLM.

\subsection{Evaluation Protocol}

We evaluate on the \texttt{ospf\_enterprise\_dhcp} scenario across all 36 fault types and three topology sizes (s, m, l), yielding 108 total cases. For the baseline and Claude-base agents, results are from complete 108-case runs. For SADE, we report results from a representative subset of 37 unique fault-size combinations that were iteratively optimized on the training topology.\footnote{Final numbers pending clean full rerun. Preliminary results are reported.}

% ============================================================
\section{Results}
% ============================================================

\subsection{Overall Comparison}

Table~\ref{tab:results} presents the overall results across all three agents.

\begin{table}[t]
\centering
\caption{Overall Agent Performance Comparison}
\label{tab:results}
\begin{tabular}{@{}lccc@{}}
\toprule
\textbf{Agent} & \textbf{Detection} & \textbf{Localization} & \textbf{RCA} \\
\midrule
ReAct/GPT-5 & 88.5\% & 56.2\% & 44.8\% \\
Claude-base & 100.0\% & 76.1\% & 67.6\% \\
SADE & \textbf{XX\%} & \textbf{XX\%} & \textbf{XX\%} \\
\bottomrule
\end{tabular}
\vspace{2pt}
{\footnotesize ReAct/GPT-5: $n=96$ valid of 108. Claude-base: $n=71$ valid of 108. SADE: $n=37$ (preliminary).}
\end{table}

\subsection{Performance by Topology Size}

Table~\ref{tab:size} shows how performance degrades with topology size across all agents.

\begin{table}[t]
\centering
\caption{RCA Accuracy by Topology Size}
\label{tab:size}
\begin{tabular}{@{}lccc@{}}
\toprule
\textbf{Agent} & \textbf{Small} & \textbf{Medium} & \textbf{Large} \\
\midrule
ReAct/GPT-5 & 43.8\% & 53.1\% & 37.5\% \\
Claude-base & 63.3\% & 74.1\% & 64.3\% \\
SADE & XX\% & XX\% & XX\% \\
\bottomrule
\end{tabular}
\end{table}

Large topologies present unique challenges: more devices to check within the same turn budget, larger diagnostic output to parse, and---as we discuss in Section~\ref{sec:selfhealing}---infrastructure instability that can cause fault evidence to disappear.

\subsection{Performance by Fault Category}

Table~\ref{tab:category} breaks down RCA accuracy by fault category, showing where SADE's layered approach provides the greatest benefit.

\begin{table}[t]
\centering
\caption{RCA Accuracy by Fault Category (Preliminary)}
\label{tab:category}
\begin{tabular}{@{}lccc@{}}
\toprule
\textbf{Category} & \textbf{ReAct} & \textbf{Claude-base} & \textbf{SADE} \\
\midrule
End-host config & XX\% & XX\% & XX\% \\
ACL blocks & XX\% & XX\% & XX\% \\
Link faults & XX\% & XX\% & XX\% \\
Traffic control & XX\% & XX\% & XX\% \\
Resource contention & XX\% & XX\% & XX\% \\
Routing & XX\% & XX\% & XX\% \\
\bottomrule
\end{tabular}
\end{table}

\subsection{Self-Healing Faults}
\label{sec:selfhealing}

Through systematic failure analysis, we identified a class of faults where the injected evidence disappears before any agent---regardless of strategy---can observe it. We term these \emph{self-healing faults} and document four instances:

\begin{enumerate}    \item \textbf{host\_incorrect\_dns (DHCP race).} The fault modifies \texttt{/etc/resolv.conf} with a wrong nameserver. However, the DHCP client (\texttt{dhclient}) with a 30-second lease overwrites the file before the agent checks, restoring correct DNS configuration.

    \item \textbf{Resource contention (OOM-kill).} Faults using \texttt{stress-ng --vm-bytes 75\%} to overload a device are OOM-killed by the host kernel under memory pressure on large topologies. The stress process disappears, leaving no observable evidence.

    \item \textbf{sender\_application\_delay (restart failure).} The fault replaces a web server's source code and restarts the service via \texttt{systemctl}. On large topologies, the service restart fails silently, leaving the original process running with normal behavior. The modified file exists on disk but the runtime behavior is unaffected.
\end{enumerate}

These self-healing faults affect \textbf{all} agents equally---they are a property of the benchmark infrastructure, not agent capability. We recommend future work on container-based network benchmarks consider fault persistence guarantees as a design requirement.

Table~\ref{tab:selfhealing} quantifies the impact. When self-healing cases are excluded from the denominator, all agents' accuracy increases, but the relative ranking is preserved.

\begin{table}[t]
\centering
\caption{Impact of Self-Healing Faults on Evaluation}
\label{tab:selfhealing}
\small
\begin{tabular}{@{}lcc@{}}
\toprule
\textbf{Agent} & \textbf{RCA (all)} & \textbf{RCA (excl. self-healing)} \\
\midrule
ReAct/GPT-5 & 44.8\% & XX\% \\
Claude-base & 67.6\% & XX\% \\
SADE & XX\% & XX\% \\
\bottomrule
\end{tabular}
\end{table}

% ============================================================
\section{Discussion}
% ============================================================

\textbf{Why layered escalation helps.}
The primary advantage of SADE over flat ReAct is \emph{systematic coverage}. The baseline agent may spend its entire turn budget investigating one hypothesis (e.g., checking resource usage on every device) and miss a simple ACL rule on a host. SADE's mandatory configuration and infrastructure sweeps in Steps 0b--0c catch low-hanging faults early, preserving budget for harder cases.

\textbf{Skill modularity vs.\ monolithic prompts.}
Our iterative development process revealed a tension: adding fault-specific rules to the main prompt improved accuracy for targeted cases but degraded it for others (prompt saturation). Moving rules to on-demand skills resolved this---the agent loads detailed knowledge only when relevant symptoms are observed.

\textbf{Container-shared kernel metrics.}
A significant pitfall in Kathara-based diagnosis is that \texttt{dmesg} and \texttt{/proc/loadavg} reflect the \emph{host} kernel state, not individual containers. Early SADE iterations that relied on \texttt{dmesg} for OOM detection produced false positives. This is a general concern for container-based network emulation benchmarks.

\textbf{Limitations.}
SADE's workflow is designed for the NIKA \texttt{ospf\_enterprise\_dhcp} topology and its 36 fault types. Generalization to other topologies and fault types requires extending the skill library. Additionally, SADE assumes a single fault per scenario; multi-fault diagnosis would require workflow modifications. The turn budget (20 turns) constrains thorough investigation on large topologies.

% ============================================================
\section{Conclusion}
% ============================================================

We presented SADE, a structured diagnostic agent architecture for LLM-based network fault diagnosis. By combining a layered escalation workflow with modular diagnostic skills, SADE achieves substantial improvements over both a ReAct/GPT-5 baseline and an unstructured Claude agent on the NIKA benchmark. Our analysis also identified self-healing faults as a previously undocumented limitation of container-based network emulation benchmarks, where fault evidence disappears before agent observation. These findings suggest that domain-structured prompting---rather than model capability alone---is the key driver of diagnostic accuracy for LLM-based network troubleshooting.

\bibliographystyle{IEEEtran}

\begin{thebibliography}{1}

\bibitem{nika}
[NIKA benchmark reference --- to be filled]

\bibitem{react}
S.~Yao, J.~Zhao, D.~Yu, N.~Du, I.~Shafran, K.~Narasimhan, and Y.~Cao, ``ReAct: Synergizing reasoning and acting in language models,'' in \emph{Proc. ICLR}, 2023.

\end{thebibliography}

\end{document}
